\section*{Abstract}
\addcontentsline{toc}{section}{Abstract}

Cities facing intensifying heat must decide where to intervene and with which
nature-based solution, typically across many more candidate sites than can be
studied in detail. The available options are unstructured judgment, which is
fast but neither comparable nor defensible, and microclimate simulation, which
is rigorous but too slow and costly for triage. This paper specifies a
screening-level, multi-criteria methodology that occupies the gap between them.
The methodology converts a structured description of an urban site and one
proposed nature-based intervention into a Heat Priority Index, a Cooling
Potential Score, an indicative daytime pedestrian-level air-temperature
reduction range, derived cooling-energy and greenhouse-gas estimates, equity
and co-benefit scores, and a single NbS Cooling Opportunity Score on a
\numrange{0}{100} scale, each accompanied by an explicit confidence rating. It
is built on a library of fourteen intervention typologies whose performance
envelopes are derived from published evidence, principally an umbrella review
of sixty-four systematic reviews \parencite{keravec2026}, an empirical
meta-analysis of park cooling \parencite{bowler2010}, and a scale-dependent
canopy study \parencite{ziter2019}; the composite-indicator construction
follows established practice \parencite{nardo2008} and the prioritisation logic
follows \textcite{norton2015}. Three calibration choices are emphasised: a
downward correction for the modelling bias present in the synthesised
literature, a refusal to claim super-additive cooling for combined
intervention packages, and the clipping of all reported temperature ranges to
the literature-supported envelope so that favourable site conditions can never
generate a claim exceeding published evidence. The methodology's limitations
are stated rather than minimised: all cooling values are daytime only; the
quantity estimated is air temperature, not thermal comfort; the evidence base
is geographically skewed towards temperate and subtropical Northern Hemisphere
cities; no default costs are shipped; and the aggregation weights are declared
expert judgment, including a disclosed equity-forward emphasis that gives
vulnerability an effective weight of \num{0.25} against heat exposure at
\num{0.15}. The sensitivity analysis specified here is planned methodology, not
a reported result. The tool, its configuration, and this document are public
under the Apache License 2.0.

\vspace{6mm}
\noindent\textbf{Keywords:} urban heat island; nature-based solutions; urban
cooling; composite indicator; screening assessment; heat vulnerability;
decision support; evidence traceability.
